% =========================================================
%  references.tex
%  Daftar pustaka buku ajar.
%  Nomor rujukan [R1]--[R35] sengaja dibuat sama dengan
%  rps/sections/10-referensi.tex agar tabel pemetaan minggu
%  dan rujukan pada RPS tetap berlaku pada buku ajar.
%  \cite{kunci} menghasilkan [Rn] sesuai nomor tersebut.
% =========================================================

\begin{sloppypar}
% Lingkungan thebibliography pada kelas book mencetak judul tanpa
% menambahkannya ke daftar isi, sehingga baris berikut diperlukan.
\addcontentsline{toc}{chapter}{\bibname}
\begin{thebibliography}{R35}
\setlength{\itemsep}{5pt}

% ---------- Referensi utama (dipakai mahasiswa) ----------
\bibitem[R1]{sweigart-mg}
  A. Sweigart, \istilah{Making Games with Python \& Pygame}, 2012.
  \url{https://inventwithpython.com/makinggames.pdf}

\bibitem[R2]{sweigart-invent}
  A. Sweigart, \istilah{Invent Your Own Computer Games with Python},
  edisi ke-4, 2018.
  \url{https://inventwithpython.com/inventwithpython.pdf}

\bibitem[R3]{pygame-docs}
  Pygame Community, \istilah{Pygame Documentation --- Quick Start dan
  Referensi API}. \url{https://www.pygame.org/docs/}

\bibitem[R4]{pygame-sprite-mask}
  Pygame Community, \istilah{Pygame Documentation --- Sprite} dan
  \istilah{Mask}. \url{https://www.pygame.org/docs/ref/sprite.html} dan
  \url{https://www.pygame.org/docs/ref/mask.html}

\bibitem[R5]{nystrom-gameloop}
  R. Nystrom, \istilah{Game Programming Patterns} --- bab
  \istilah{Game Loop}.
  \url{https://gameprogrammingpatterns.com/game-loop.html}

\bibitem[R6]{fiedler-timestep}
  G. Fiedler, ``Fix Your Timestep!'', \istilah{Gaffer On Games}, 2004.
  \url{https://gafferongames.com/post/fix_your_timestep/}

\bibitem[R7]{catto-integration}
  E. Catto, \istilah{Numerical Integration}, GDC 2009.
  \url{https://box2d.org/files/ErinCatto_NumericalIntegration_GDC2009.pdf}

\bibitem[R8]{santos-pathfinding}
  E. dos Santos, \istilah{Game AI: Pathfinding}, bahan kuliah, 2020.
  \url{https://edirlei.com/aulas/game-ai-2020/GAME_AI_Lecture_02_Pathfinding_2020.pdf}

\bibitem[R9]{santos-fsm}
  E. dos Santos, \istilah{Game AI: Finite State Machines}, bahan kuliah,
  2020.
  \url{https://edirlei.com/aulas/game-ai-2020/GAME_AI_Lecture_03_Finite_State_Machines_2020.pdf}

\bibitem[R10]{bendersky-pygame-astar}
  E. Bendersky, ``Writing a Game in Python with Pygame, Part III:
  Pathfinding'', 2009.
  \url{https://eli.thegreenplace.net/2009/01/09/writing-a-game-in-python-with-pygame-part-iii}

\bibitem[R11]{munir-astar}
  R. Munir, \istilah{Penentuan Rute, Bagian 2: Algoritma A*}, bahan
  kuliah Institut Teknologi Bandung.
  \url{https://informatika.stei.itb.ac.id/~rinaldi.munir/Stmik/2020-2021/Route-Planning-Bagian2-2021.pdf}

\bibitem[R12]{unity-leveldesign}
  Unity Technologies, \istilah{Introduction to Level Design in Game
  Development and in Unity}, e-book resmi.
  \url{https://assets2.brandfolder.io/bf-boulder-prod/khmctn58vjqkr42wmm3mqvq9/v/1131033155/original/UNITY_Ebook_Introduction%20to%20level%20design%20in%20game%20development%20and%20in%20Unity_V2.pdf}

\bibitem[R13]{unity-learn}
  Unity Technologies, \istilah{Unity Learn --- Unity Essentials} dan
  \istilah{Junior Programmer: Create with Code}.
  \url{https://learn.unity.com/}

\bibitem[R14]{unity-manual-import}
  Unity Technologies, \istilah{Unity Manual --- Importing Assets} dan
  \istilah{Audio}.
  \url{https://docs.unity3d.com/Manual/ImportingAssets.html}

% ---------- Referensi pendukung: RPS dan materi perguruan tinggi ----------
\bibitem[R15]{its-rps-game}
  Institut Teknologi Sepuluh Nopember, \istilah{RPS Desain dan
  Pemrograman Game}, Program Studi Teknik Komputer.
  \url{https://www.its.ac.id/komputer/wp-content/uploads/sites/28/2018/03/Desain-dan-pemograman-Game.pdf}

\bibitem[R16]{unima-rps-game}
  Universitas Negeri Manado, \istilah{RPS Pengembangan Game}, Program
  Studi Teknik Informatika, 2025/2026.
  \url{https://ti.unima.ac.id/rps/view/11}

\bibitem[R17]{unesa-rps-game}
  Universitas Negeri Surabaya, \istilah{RPS Pemrograman Game}, S1
  Pendidikan Teknologi Informasi.
  \url{https://sidia.unesa.ac.id/rps-pdf/s1-pendidikan-teknologi-informasi/pemrograman-game.pdf}

\bibitem[R18]{polibatam-kurikulum}
  Politeknik Negeri Batam, \istilah{Kurikulum Program Studi Teknologi
  Permainan}, 2024.
  \url{https://www.polibatam.ac.id/wp-content/uploads/2024/06/Kurikulum-Prodi-TP.pdf}

\bibitem[R19]{univrab-pemrograman-game}
  Universitas Abdurrab, \istilah{Deskripsi Mata Kuliah Pemrograman Game},
  SPADA. \url{http://spada.univrab.ac.id/course/info.php?id=1303&lang=en}

\bibitem[R20]{uma-pygame}
  Universitas Medan Area, \istilah{Panduan Membuat Game 2D Sederhana
  dengan Pygame}, 2025.
  \url{https://p2dpm.uma.ac.id/2025/10/15/panduan-membuat-game-2d-sederhana-dengan-pygame/}

\bibitem[R21]{ub-pygame}
  Universitas Brawijaya, ``Game Development with Python: Let's Build a
  Simple Game!'', \istilah{Gramaswara}.
  \url{https://gramaswara.ub.ac.id/index.php/gramaswara/article/download/185/128}

\bibitem[R22]{dinus-pygame}
  Universitas Dian Nuswantoro, ``Pemrograman dan Game dengan
  Python/Pygame'', \istilah{Abdimasku}.
  \url{https://abdimasku.lppm.dinus.ac.id/index.php/jurnalabdimasku/article/download/2965/735}

\bibitem[R23]{itb-collision}
  Institut Teknologi Bandung, ``Collision Detection in Precision
  Parry-Based Combat'', makalah IF2123, 2024.
  \url{https://informatika.stei.itb.ac.id/~rinaldi.munir/AljabarGeometri/2024-2025/Makalah/Makalah-IF2123-Algeo-2024%20(31).pdf}

\bibitem[R24]{stmikjayakarta-pygame}
  STMIK Jayakarta, ``Penerapan Library Pygame dalam Game RPG The
  Adventure'', \istilah{Jurnal JMI Jayakarta}.
  \url{https://journal.stmikjayakarta.ac.id/index.php/JMIJayakarta/article/download/1013/677/}

\bibitem[R25]{itn-astar-fsm}
  Institut Teknologi Nasional Malang, ``Penerapan A* Pathfinding dan
  Finite State Machine pada Game Lost Civilization'', \istilah{JATI}.
  \url{https://ejournal.itn.ac.id/index.php/jati/article/download/5402/3593}

\bibitem[R26]{umri-astar-behaviour}
  Universitas Muhammadiyah Riau, ``Penerapan Algoritma A* dan Behaviour
  Trees untuk NPC'', \istilah{COSCITECH}.
  \url{https://ejurnal.umri.ac.id/index.php/coscitech/article/download/5418/2453/17612}

\bibitem[R27]{uinmalang-fsm-fuzzy}
  UIN Maulana Malik Ibrahim Malang, \istilah{Implementasi FSM dan Logika
  Fuzzy untuk Perilaku NPC}, skripsi.
  \url{http://etheses.uin-malang.ac.id/8204/1/11650053.pdf}

\bibitem[R28]{uny-katelu}
  Universitas Negeri Yogyakarta, \istilah{Pengembangan Permainan Edukasi
  KATELU dengan Unity 3D}, skripsi.
  \url{https://eprints.uny.ac.id/52540/1/TAS_AFIF%20FATCHUR%20RAHMAN_12520241062.pdf}

\bibitem[R29]{undip-media}
  Universitas Diponegoro, ``Media Pembelajaran Budaya Indonesia
  Menggunakan Unity'', \istilah{Jurnal Teknologi dan Sistem Komputer}.
  \url{https://jtsiskom.undip.ac.id/article/download/12667/12286}

\bibitem[R30]{untan-edukasi}
  Universitas Tanjungpura, ``Game Edukasi Budaya Indonesia'',
  \istilah{JCSKom}.
  \url{https://jurnal.untan.ac.id/index.php/jcskommipa/article/download/49393/75676592053}

\bibitem[R31]{rokania-unity}
  STKIP Rokania, \istilah{Pelatihan Dasar Pemrograman Game dengan Unity},
  \istilah{Jurnal MNR}.
  \url{https://e-jurnal.rokania.ac.id/index.php/jmnr/article/download/39/28/172}

\bibitem[R32]{pygame-chimp}
  Pygame Community, \istilah{Pygame Tutorial: Line-by-Line Chimp}.
  \url{https://www.pygame.org/docs/tut/ChimpLineByLine.html}

\bibitem[R33]{reintech-game-ai}
  Reintech, ``How to Create a Video Game AI with Python''.
  \url{https://reintech.io/blog/how-to-create-a-video-game-ai-with-python}

\bibitem[R34]{itb-astar-navigasi}
  Institut Teknologi Bandung, ``Implementasi Algoritma A* dalam Navigasi
  Karakter'', makalah IF2120, 2023.
  \url{https://informatika.stei.itb.ac.id/~rinaldi.munir/Matdis/2022-2023/Makalah2022/Makalah-Matdis-2022%20(86).pdf}

\bibitem[R35]{stmikdharmapala-astar}
  STMIK Dharmapala Riau, ``Rancang Bangun Pathfinding A* pada NPC
  Menggunakan Unity Engine'', \istilah{Jurnal JIKB}.
  \url{https://ojs.stmikdharmapalariau.ac.id/index.php/jikb/article/download/831/574}

\end{thebibliography}
\end{sloppypar}
